\subsection{Canonical log structures on prestable curves}%
\label{sub:Canonical log structures on prestable curves}

Let \(\mf{M}_{g,n}\) be the moduli of \(n\)-pointed prestable curves. Then the singular locus is a normal crossings divisor, so there is a canonical log structure
\[ M_{\mf{M}_{g,n}}^{\on{can}} = M_{\partial \mf{M}_{g,n}}. \]
If we consider the universal family \(\mf{C}_{g,n}\), we can form \(\pi^{-1}(\partial \mf{M}_{g,n}) \subseteq \mf{C}_{g,n}\). Therefore, we have a map
\[ \pi^* M_{\partial \mf{M}_{g,n}} \to M_{\pi^{-1}(\partial \mf{M}_{g,n})}. \]
We can also consider the divisors \(p_i(\mf{M}_{g,n}) \subset \mf{C}_{g,n}\), which also define log structures
\[ P_i \coloneqq M_{p_i(\mf{M}_{g,n})}. \]
Therefore, we can define a log structure
\[ M_{\mf{C}_{g,n}}^{\on{can}} \coloneqq M_{\pi^{-1}(\partial\mf{M}_{g,n})} \oplus_{\mc{O}^{\times}} \bigoplus_{i=1}^n P_i. \]
Then if we consider
\[ \pi^{\flat} \colon \ul{\pi}^{-1} M_{\mf{M}_{g,n}}^{\on{can}} \to M_{\pi^{-1}(\partial \mf{M}_{g,n})} \to M_{\mf{C}_{g,n}}^{\on{can}}, \]
we can lift \(\pi\) to a map of log stacks.

\begin{thm}[{\cite{FKato2000}}]
    The map \(\pi\) is a log smooth integral map.
\end{thm}

\begin{defn}
    For a family \((\ul{C}/\ul{S}, \vec{p})\) of prestable curves, we can pull back along the diagram
    \begin{equation*}
    \begin{tikzcd}
        \ul{C} \ar{r} \ar{d} & \mf{C}_{g,n} \ar{d} \\
        S \ar{r} & \mf{M}_{g,n}
    \end{tikzcd}
    \end{equation*}
    and endow the left edge with a canonical log structure.
\end{defn}

\begin{rmk}
    Let \(\ul{S} = \Spec k\). If \(S^{\on{can}} = \Spec (\N^r \to k)\), then \(r\) is the number of nodes on \(\ul{C}\). Working \'etale locally on a node \(q\), the log structure looks like
    \begin{equation*}
    \begin{tikzcd}
        \N^{r-1} \oplus \N^2 \ar{r} & k[x,y]/xy \\
        \N^{r-1} \oplus \N \ar{u}{\on{id}\oplus\Delta} \ar{r} & k \ar{u}.
    \end{tikzcd}
    \end{equation*}
\end{rmk}

\begin{thm}[{\cite{FKato2000}}]
    Let \((\mc{C}/S, \vec{p})\) be an \(n\)-pointed log curve on an fs base \(S\). Then there exists a unique \(M_S^{\on{can}} \to M_S\) such that
    \begin{equation*}
    \begin{tikzcd}
        \mc{C} \ar{d} \ar{r} & \mc{C}^{\on{can}} \ar{d} \\
        S \ar{r}{a} & S^{\on{can}}
    \end{tikzcd}
    \end{equation*}
    is a Cartesian diagram in the category of fs log schemes.
\end{thm}

A slogan is that a log curve is a prestable curve together with \(M_S^{\on{can}} \to M_S\).

\begin{exm}
    Let \(\ul{S} = \Spec k[s]\) and \(\ul{T} = \Spec k[t]\). Also choose \(\ell \in \N\). Then consider the diagram
    \begin{equation*}
    \begin{tikzcd}
        \Spec \frac{k{ [x,y,s] }}{xy-s^{\ell}} \ar{d} \ar{r} & \Spec \frac{k{ [x,y,s] }}{xy-s^{\ell}} \ar{d} \ar{r} & \Spec  \frac{k{[x,y,t]}}{xy-t} \ar{r}\ar{d} & \mf{C}_{g,n} \ar{d} \\
        \Spec k{[s]} \ar{r}{\ul{ \on{id} }} & \Spec {k[s]} \ar{r} & \Spec k{[t]} \ar{r} & \mf{M}_{g,n}.
    \end{tikzcd}
    \end{equation*}
    On log structures, consider the case where the bottom left arrow corresponds to \(1 \mapsto \ell\) on \(\N\). Then at the top left corner, the monoid becomes \(\N^2 \oplus_{\N, \ell} \N\). If we call this \(P\), then the map \(k[P] \to k[\N]\) is not a family of log curves with the canonical log structure, but is a family of log curves.
\end{exm}

\subsection{Moduli of log curves}%
\label{sub:Moduli of log curves}

\begin{defn}
    A map between \(n\)-pointed log curves \((\mc{C}, S) \to (\mc{C}', S')\) is a pair \((\ell \colon \mc{C} \to \mc{C}', \Theta \colon S \to S')\) such that \(\Theta\) is strict and the diagram
    \begin{equation*}
    \begin{tikzcd}
        \mc{C} \ar{d} \ar{r}{\ell} & \mc{C}' \ar{d} \\
        S \ar{r}{\Theta} & S'
    \end{tikzcd}
    \end{equation*}
    is Cartesian. Denote by \(\mf{M}_{g,n}^{\log}\) the fibered category parameterizing \(n\)-pointed log curves.
\end{defn}

\begin{prop}[{\cite{FKato2000,Olsson2003}}]
    We have \(\mf{M}_{g,n}^{\log} \simeq \ms{Log}_{(\mf{M}_{g,n}, M_{\mf{M}_{g,n}}^{\on{can}})}\).
\end{prop}

In particular, for a scheme \(\ul{T}\), the groupoid \(\mf{M}_{g,n}^{\log}(\ul{T})\) parameterizes log structures \(M_T\) together with log curves \(\mc{C} / T\).

Consider a log curve \(\pi \colon \mc{C} \to S, \vec{p}\) over \(S = \Spec (Q \to k)\). Denote \(\bar{\pi}^{\flat} \colon \bar{\pi}^* \bar{M}_S \to \bar{M}_{\mc{C}}\). At a generic point or a general smooth point \(\eta\), the map is
\[ \bar{\pi}_{\eta}^{\flat} \colon (\pi^{-1} \bar{M}_S)_{\eta} = Q \simeq \bar{M}_{\mc{C}, \eta}. \]
At a marked point \(p_i \in \ol{\ab\{\eta\}}\), we have \(\bar{M}_{\mc{C}, p_i} = Q \oplus \N\). In this case we have
\[ \chi_{\eta} \colon \bar{M}_{\mc{C}, p_i} \xrightarrow{\on{pr}_1} \bar{M}_{\mc{C}, \eta} \]
the specialization map and the pullback
\[ \bar{\pi}_{p_i}^{\flat} \colon Q \hookrightarrow Q \oplus \N. \]
At a node \(q\), we have \(\bar{M}_{\mc{C}, q} = Q \oplus_{\N} \N^2\). Let \(\rho_q\) denote the image of \(1\) in \(Q\), which is the edge length or smoothing parameter at \(q\). Denote by \(\bar{\eta}_1\) and \(\bar{\eta}_2\) the generic points of the two branches of \(q\). Then the generalization maps are
\begin{align*}
    \chi_{\eta_1} \colon Q \oplus_{\N} \N^2 \to Q \qquad (q, (a,b)) \mapsto q + a \rho_q ;\\
    \chi_{\eta_2} \colon Q \oplus_{\N} \N^2 \to Q \qquad (q, (a,b)) \mapsto q + b \rho_q .
\end{align*}

\subsection{Tropicalization of log curves}%
\label{sub:Tropicalization of log curves}

We will only consider the situation over a log point \(S = \Spec(Q \to k)\). We will consider the category of cones \(\ms{Cone}\), whose objects are pairs \((w, N_w)\) such that \(w \subseteq N_w \otimes_{\Z} \R\), where \(w\) is a full-dimensional strongly convex rational polyhedral cone. The morphisms are maps \((w_1, N_{w_1}) \to (w_2, N_{w_2})\), which consist of
\[ \chi \colon N_{w_1} \to N_{w_2} \]
such that \(\chi_{\R}(w_1) \subseteq w_2\). 

\begin{exm}
    Let \(X\) be an fs log scheme and consider \(( (\bar{M}_{X,x}^{\vee})_{\R}, (\bar{M}_{X,x}^{\vee})_{\Z} )\).
\end{exm}

We will now consider the category \(\ms{GCC}\) of generalized cone complexes. Its objects are topological spaces \(\abs{\Sigma}\) together with a presentation as the colimit of a diagram in \(\ms{Cone}\) with all arrows being inclusions of faces. Its morphisms are continuous map \(\abs{\Sigma} \to \abs{\Sigma'}\) such that for each \((w, N_w) \in \Sigma\), the map \(w \to \abs{\Sigma'}\) factors through \((w, N_w) \to (w', N_{w'})\).

\begin{exm}
    If \(f \colon X \to Y\) is a morphism of fs log schemes, the map \(\bar{f}^{\flat} \colon f^* \bar{M}_Y \to \bar{M}_X\) induces \(\Sigma(X) \to \Sigma(Y)\) on tropicalizations.
\end{exm}

\begin{defn}
    A \textit{cone complex} is a generalized cone complex \(\Sigma\) such that for each \((\sigma, N_{\sigma}) \in \Sigma\), the map \(\sigma \to \abs{\Sigma}\) is injective.
\end{defn}

\begin{thm}[{\cite{ACP2015}}]
    If \(X\) is fs and log smooth over \(\ul{\C}\) with Zariski log structure, then \(\Sigma(X)\) is a cone complex.
\end{thm}

For a monoid \(Q\), we will denote \(w_Q = Q_{\R}^{\vee} = \Hom(Q, \R_{\geq 0})\) and the lattice \(\Hom(Q, \Z) = Q_{\Z}^{\vee}\). At a smooth point, we have
\[ w_{\eta} = (Q_{\R}^{\vee}, Q_{\Z}^{\vee}). \]
At a marked point, we see that
\[ w_p = \ab(\Hom (Q \oplus \N, \R_{\geq 0}) = Q_{\R}^{\vee} \oplus \R_{\geq 0}, Q_{\Z}^{\vee}\oplus \Z). \]
At a node, we have
\begin{align*}
    w_q &= \ab(\Hom(Q \oplus_{\N} \N^2, \R_{\geq 0}), ( \bar{M}_{\mc{C}, q} )_{\Z}^{\vee}).
\end{align*}
Note that the cone is in fact equal to
\[ \ab\{(m, m_x, m_y) \mid m(\rho_q) = m_x + m_y \} \simeq \ab\{(m, e_x) \mid e_x \leq m(\rho_q) \} \subseteq Q_{\R}^{\vee} \times \R_{\geq 0}. \]

Now that we understand the cones associated to every point, we now need to discuss the face maps. The generalization maps corresponding to marked points are
\[ Q \oplus \N \xrightarrow{\on{pr}_1} Q, \]
and the associated face maps are 
\[ w_{\eta} \hookrightarrow w_p \qquad m \mapsto (m,0). \]
At a node, the two face maps are
\[ \chi_{\eta_1}^{\vee} \colon Q_{\R}^{\vee} \to w_q \qquad m \mapsto (m, m(\rho_q)) \]
and 
\[ \chi_{\eta_2}^{\vee} \colon Q_{\R}^{\vee} \to w_q \qquad m \mapsto (m, 0). \]

\begin{exm}
    Consider a curve which is a union of two rational curves with a marked point \(p\) on one branch. Then let \(\ul{S} = \Spec \C\). Then we see that
    \[ S^{\on{can}} = \Spec(\N \to \C), \]
    \(Q = \N\), and \(Q_{\R}^{\vee} = \R_{\geq 0}\). We also see that \(\bar{M}_{C, q} = \N \oplus_{\N} \N^2\) and \(\bar{M}_{C, p} = \N \oplus \N\).
\end{exm}

\begin{figure}[H]
    \centering
    \begin{tikzpicture}[
        x=1cm,
        y=1cm,
        cone ray/.style={thick,->,>=stealth},
        common face/.style={red,line width=2pt,line cap=round},
        diagonal face/.style={blue!75!black,line width=2pt,line cap=round}
    ]
        % The cone at the marked point.
        \coordinate (p) at (0,4.45);
        \draw[cone ray] (p) -- ++(-0.55,2.15);
        \draw[cone ray] (p) -- ++(2.25,1.85);
        \draw[diagonal face]
            ($(p)+(0.75,0.62)$) -- ($(p)+(1.45,1.19)$);
        \node[above] at ($(p)+(0.55,2.2)$)
            {$w_p=\R_{\geq 0}\times\R_{\geq 0}$};

        % The cone at the first generic point.
        \coordinate (etaone) at (0.35,4);
        \draw[cone ray] (etaone) -- ++(2.25,1.85);
        \draw[diagonal face]
            ($(etaone)+(0.75,0.62)$) -- ($(etaone)+(1.45,1.19)$);
        \node[above right] at ($(etaone)+(2.0,1.6)$) {$w_{\eta_1}$};

        % The cone at the node.
        \coordinate (q) at (0.35,3.25);
        \draw[cone ray] (q) -- ++(2.25,1.85);
        \draw[cone ray] (q) -- ++(3.35,0);
        \draw[diagonal face]
            ($(q)+(0.75,0.62)$) -- ($(q)+(1.45,1.19)$);
        \draw[common face] ($(q)+(0.85,0)$) -- ($(q)+(1.75,0)$);
        \node at ($(q)+(2.55,0.75)$) {$w_q$};
        \node[anchor=west] at ($(q)+(4.05,0.75)$)
            {$\bigl\{(a,b)\mid b\leq a\bigr\}\subseteq\R_{\geq 0}^{2}$};

        % The colored intervals indicate the common faces used in the gluing.
        \coordinate (etatwo) at (0.35,2.75);
        \draw[cone ray] (etatwo) -- ++(3.35,0);
        \draw[common face] ($(etatwo)+(0.85,0)$) -- ($(etatwo)+(1.75,0)$);
        \node[right] at ($(etatwo)+(3.45,0)$) {$w_{\eta_2}$};
    \end{tikzpicture}
\end{figure}


\subsection{Tropical curves}%
\label{sub:Tropical curves}

\begin{defn}
    A family of \textit{tropical curves} \((G, \vec{g}, \vec{\ell})\) over a cone \(w, N_w\) consists of
    \begin{enumerate}
        \item A connected graph \(G = V(G) \cup E(G) \cup L(G)\);
        \item A leg ordering \(\vec{m} \colon L(G) \xrightarrow{\sim} \ab\{1, \ldots, n\}\);
        \item A genus function \(\vec{g} \colon V(G) \to \N\);
        \item An edge length function \(\ell \colon E(G) \to \Hom(w \cap N_w, \N) \setminus \ab\{0\}\).
    \end{enumerate}
\end{defn}

Now we will construct \(\Sigma(G, \ell)\). For \(v \in V(G)\), we set \(w_v = w\). For \(e \in E(G)\), we set
\[ w_e = \ab\{ (s, \lambda) \mid \lambda \leq \vec{\ell}(e)(s) \} \subseteq w \times \R_{\geq 0}. \]
For legs \(\ell \in L(G)\), we set \(w_{\ell} = w \times \R_{\geq 0}\). We now need to define the face maps. If \(\ell\) is incident to \(v\), then the face map is
\[ w_v \hookrightarrow w_{\ell} \colon s \mapsto (s, 0). \]
If \(e\) is incident to two vertices \(v_1, v_2\), then the face maps are
\begin{align*}
    w_{v_1} \hookrightarrow w_e &\colon s \mapsto (s, \vec{\ell}(e)(s)) \\
    w_{v_2} \hookrightarrow w_e & \colon s \mapsto (s, 0).
\end{align*}
\begin{rmk}
    The choice of \(v_1\) and \(v_2\) does not change the cone complex up to isomorphism, but it does matter when we want to define contact orders for log maps. However, in general there is no canonical way to make this choice.
\end{rmk}

\subsection{Punctured curves}%
\label{sub:Punctured curves}

Punctured logarithmic curves and maps are developed in~\cite{ACGS2025}.

We consider a nodal curve
\[ \Spec k[x,y]/xy \to \Spec k \]
with log structures given by \(\N^2\) with chart \((a,b) \mapsto x^a y^b\) and \(\N \to k\) being the standard log point. Now consider \(\ul{C}_1 = Z(x) = \Spec k[y]\) and mark the node by \(p \colon \Spec k \to \Spec k[y]\). Note that
\[ C^{\on{can}} = (\ul{C}, M_C) \to S^{\on{can}} = \Spec(\N \to k). \]
Denote \(M_{C_1}^{\circ} = M_C|_{C_1}\). Then we see that on log structures we have the diagram
\begin{equation*}
\begin{tikzcd}
    k[y] & & \N^2 \ar{ll}[swap]{(a,b) \mapsto 0^a y^b} \\
    k \ar{u} & & \N \ar{ll} \ar{u}{\Delta}.
\end{tikzcd}
\end{equation*}
In particular, we can compute the stalk \(\bar{M}_{C_1}^{\circ}|_p = \N^2\), and therefore we have
\[ \bar{\pi}_{C_1^{\circ}}^{\flat} |_p \colon \N \xrightarrow{\Delta} \N^2. \]
However, we should have the map
\[ \bar{\pi}^{\flat}_{C_1}|_p \colon \N \xrightarrow{i_1} \N \oplus \N, \]
and therefore \(C_1\) is \textbf{not} a log curve. We can then consider the diagram
\begin{equation*}
\begin{tikzcd}
    \N \oplus \Z \\
    \N^2 \ar[hookrightarrow]{u} \ar{rr}{(0,1)\mapsto (0,1)}[swap]{(1,0)\mapsto (1,1)} & & \N^2 \ar{ull}[swap]{\Vectorstack{{(1,0)\mapsto(1,-1)} {(0,1)\mapsto(0,1)}}} \\
    & \N \ar{ul}{i_1} \ar{ur}{\Delta}.
\end{tikzcd}
\end{equation*}
Therefore, we will consider extending the log structures.

Recall that for a log curve, we have \(M_{\mc{C}} = M_{\mc{C}}' \oplus_{\mc{O}_{\mc{C}}^{\times}} P_{\mc{C}}\), where \(P_{\mc{C}}\) is the divisorial log structure by the markings.

\begin{defn}
    A \textit{punctured curve} over an fs log scheme \(S\) consists of \((\mc{C}^{\circ} \xrightarrow{p} \mc{C} \xrightarrow{\pi} S, \{p_i\}_{i=1}^n)\), where \((\mc{C} \xrightarrow{\pi} S, \vec{p})\) is a log curve and the map \(p \colon \mc{C}^{\circ} \to \mc{C}\) is a map of fine log schemes satisfying the following conditions:
    \begin{enumerate}
        \item The underlying map \(\ul{p}\) is an isomorphism;
        \item On log structures, the map \(p^{\flat} \colon M_{\mc{C}} \hookrightarrow M_{\mc{C}^{\circ}}\) is an isomorphism away from the \(p_i\) and induces
            \[ M_{\mc{C}} \xrightarrow{p^{\flat}} M_{\mc{C}^{\circ}} \hookrightarrow M_{\mc{C}} \oplus P^{\mr{gp}}. \]
        \item For any marking \(p_i\) and any \(s \in M_{\mc{C}^{\circ}, p_i} \setminus M_{\mc{C}, p_i}\), we have \(\alpha_{\mc{C}^{\circ}}(s) = 0\).
    \end{enumerate}
\end{defn}

Roughly speaking, a punctured curve is a log curve plus the data of a commutative diagram
\begin{equation*}
\begin{tikzcd}
    M_{\mc{C}^{\circ}} \ar{dr} \ar[hookrightarrow]{r} & M_{\mc{C}}' \oplus_{\mc{O}_{\mc{C}}^{\times}} P^{\on{gp}} \\
    & \mc{O}_{\mc{C}} \\
    M_{\mc{C}} \ar[hookrightarrow]{uu} \ar{ur}.
\end{tikzcd}
\end{equation*}


\begin{exm}
    Let \((\ul{\mc{C}} \to \ul{S} = \Spec k, p)\) be a smooth \(1\)-pointed curve.
    \begin{enumerate}
        \item Let \((\mc{C}/S, p)\) be endowed with the canonical log structure (so \(S\) has no log structure). Let \(M_{\mc{C}} = P\). Then any puncturing \(p^{\flat} \colon M_{\mc{C}} \hookrightarrow M_{\mc{C}^{\circ}} \subset P^{\on{gp}}\) must be trivial. The reason is that \(M_{\mc{C}^{\circ}}\) cannot have a section \(s\) such that \(\alpha_{\mc{C}^{\circ}}(s) = 0\). If \(s \in P = M_{\mc{C}}\), then \(\alpha_{\mc{C}^{\circ}}(s) = \alpha_{\mc{C}}(s)\). Otherwise, if \(s \notin P\), then because \(\alpha_{\mc{C}^{\circ}}(-s)\) is nonzero, so is \(\alpha_{\mc{C}^{\circ}}(s)\).
        \item Let \(M_{\mc{C}^{\circ}} \subseteq (\N \oplus \mc{O}_{\mc{C}}^{\times}) \oplus P^{\on{gp}}\) be a puncturing. Choose a local generator \(x\)  of \(\mc{O}(-p)\). Then any \(s \in M_{\mc{C}^{\circ}, p}\) has the form \(s = ((n, \varphi), x^m)\) where \(n \in \N\) and \(\varphi \in \mc{O}_{\mc{C}, p}^{\times}\). If \(m < 0\), then we require that \(\alpha_{\mc{C}^{\circ}}(s) = 0\). This implies that \(\alpha_{\mc{C}}((n, \varphi), 0) = 0\), and therefore
            \[ \bar{M}_{\mc{C}^{\circ}, p} \subseteq \ab\{(n, m) \in \N \oplus \Z \mid m \geq 0 \text{ if } n = 0 \}. \]
            In fact, any fine submonoid of the right hand side containing \(\N \oplus \N\) can be realized by a puncturing.
        \item Let \(S = \Spec(\N \to k)\). Consider \(M_{\mc{C}^{\circ}}\) such that \(\bar{M}_{\mc{C}^{\circ}, p} = \N^2 + \N (2,-2)\). Because \((1,-1) \notin \bar{M}_{\mc{C}^{\circ}, p}\), we see that \(M_{\mc{C}^{\circ}}\) is not saturated.
    \end{enumerate}
\end{exm}

\begin{lem}[{\cite{ACGS2025}}]
    Let \(Q\) be a sharp fs monoid. Let \(w \coloneqq \Hom(Q, \R_{\geq 0}) = Q_{\R}^{\vee}\). Let \(Q^{\circ} \subseteq Q \oplus \Z\) be finitely generated with \(Q \oplus \N \subsetneq Q^{\circ}\) and \(Q^{\circ} \cap \ab\{0\} \times \Z_{<0} = \emptyset\). Then there exists a nonzero, concave, piecewise linear function
    \[ \ell \colon w \to \R_{\geq 0} \]
    with rational slopes such that
    \[ (Q^{\circ})_{\R}^{\vee} = \ab\{ (s, \lambda) \in w \times \R_{\geq 0} \mid 0 \leq \lambda \leq \ell(s) \}. \]
\end{lem}

\begin{exm}
    Let \(Q = \N^2\) and set 
    \begin{align*}
        Q^{\circ} &= \ab<((1,0), -1), ((0,1), -1), ((0,0),1)> \subset \N^2 \oplus \Z \\
        &= \ab\{((a,b), m) \mid a+b+m \geq 0, a,b \geq 0 \}.
    \end{align*}
    Then
    \[ (Q^{\circ})_{\R}^{\vee} = \ab\{ (e_x, e_y, e_t) \mid 0 \leq e_t \leq \min(e_x, e_y) \}. \]
\end{exm}





